\chapter{Polygamma row sums at complex multiplication points}
\label{ch:06-cm}

Rational gamma grids account for periods that reappear in modular forms.
At imaginary quadratic arguments the polygamma series becomes a row of a
lattice Eisenstein sum. Symmetry isolates elementary components; modular
periods and class-field coordinates account for the non-elementary ones.
Individual values and full row sums require different evidence.

\section{Lattice normalization and row sums}
For even $w\ge4$ and $\Im\tau>0$, define the absolutely convergent
\[
G_w(\tau)=\sum_{(m,n)\in\mathbb Z^2\setminus\{(0,0)\}}
\frac1{(m+n\tau)^w}=2\zeta(w)E_w(\tau).
\]
Pair positive and negative $n$, and use the defining polygamma series:
\begin{equation}\label{cm:eq:rowsum}
G_w(\tau)=2\zeta(w)+\frac2{(w-1)!}
\sum_{n\ge1}\left[\psi^{(w-1)}(n\tau)+\psi^{(w-1)}(1-n\tau)\right].
\end{equation}
The hexagonal point in this chapter is $\rho=e^{i\pi/3}$.
The modular invariant $\mathfrak g_2=E_4/\eta^8$ below is unrelated to the
second Stieltjes constant. Algebraic fits and historical numerical
residuals are labeled as such; class polynomials are formed from normalized
modular functions, rather than untreated modular-form values.

\section{The elementary layer: reflection composed with conjugation}\label{cm:sec:elementary}

Before any CM input, one family of components is fixed by symmetry alone. The Schwarz
reflection $\psi^{(m)}(\bar z)=\overline{\psi^{(m)}(z)}$ turns the reflection law into a
statement about a single real component, leaving the other component untouched.

\begin{result}[vertical-line trigamma law]\label{cm:res:line}
On the line $\operatorname{Re}z=\tfrac12$ the reflection law alone evaluates the real part of
the trigamma function at every height: for all real $t$,
\[
\operatorname{Re}\psi^{(1)}\!\Bigl(\tfrac12+it\Bigr)=\frac{\pi^2}{2\cosh^2(\pi t)} .
\]
In particular $\operatorname{Re}\psi^{(1)}\!\bigl(\tfrac{1+i}{2}\bigr)=\pi^2/(2\cosh^2(\pi/2))$
\dps{residual $0$ at $50$ dps}. Off the symmetric line, combining reflection, conjugation and
the recurrence gives e.g.
\[
\operatorname{Re}\psi^{(1)}(i)=-\frac12-\frac{\pi^2}{2\sinh^2\pi}
\qquad\text{\dps{residual $0$ at $50$ dps}.}
\]
The proof of the first identity is one line: on $\operatorname{Re}z=\tfrac12$ one has
$1-z=\bar z$, so the reflection law
$\psi^{(1)}(z)+\psi^{(1)}(1-z)=\pi^2/\sin^2(\pi z)$ reads
$2\operatorname{Re}\psi^{(1)}(z)=\pi^2/\sin^2(\pi z)$ there.
\end{result}

The imaginary parts on this line stay opaque --- they are odd-row lattice data in the sense of
Section~\ref{cm:sec:opaque}.

\subsection{The FoxTrot identity demystified}

MathWorld's \emph{Digamma Function} page \cite{cm:MathWorldDigamma} calls the following
identity ``surprising''; it arises from the FoxTrot series $\sum(-1)^{n-1}n^2/(n^3+1)$,
whose partial fractions hit the roots of $n^2-n+1$:
\begin{equation}\label{cm:eq:foxtrot}
-\psi\bigl(\tfrac{\omega}{2}\bigr)-\psi\bigl(-\tfrac{\omega^2}{2}\bigr)
+\psi\bigl(\tfrac{1+\omega}{2}\bigr)+\psi\bigl(\tfrac{1-\omega^2}{2}\bigr)
=2\pi\,\sech\Bigl(\tfrac{\sqrt3\,\pi}{2}\Bigr),\qquad \omega=e^{i\pi/3}.
\end{equation}

\begin{law}[line law for digamma differences]\label{cm:law:lineline}
For every shift $c\in(0,1)$, on the line $\operatorname{Re}z=(1-c)/2$,
\[
\operatorname{Re}\bigl[\psi(z+c)-\psi(z)\bigr]=\pi\,\operatorname{Re}\cot(\pi\bar z).
\]
\dps{verified for $c=\tfrac12,\tfrac13,\tfrac23$ at generic heights $t$}
\end{law}

The mechanism is the same one-liner: on that line the shift satisfies $z+c=1-\bar z$, so
reflection composed with conjugation eliminates the digamma values, leaving an elementary
$\sinh/\cos$ closed form. The FoxTrot identity \eqref{cm:eq:foxtrot} is the instance
$c=\tfrac12$, $z=\tfrac14+i\tfrac{\sqrt3}{4}$: its four arguments pair into
$2\operatorname{Re}[\psi(z+\tfrac12)-\psi(z)]$, and on $\operatorname{Re}z=\tfrac14$ one has
$\pi\operatorname{Re}\cot(\pi\bar z)=\pi\sech(2\pi t)$. The identity is thus the point
$t=\sqrt3/4$ on a line of identities, itself one member of a two-parameter family --- no
complex multiplication required.

Differentiating in $t$ ladders the law up the polygamma tower
\dps{weights $2,3$ verified at $10^{-50}$}:
\begin{align*}
\operatorname{Im}\bigl[\psi^{(1)}(\tfrac34+it)-\psi^{(1)}(\tfrac14+it)\bigr]
&=2\pi^2\sech(2\pi t)\tanh(2\pi t),\\
\operatorname{Re}\bigl[\psi^{(2)}(\tfrac34+it)-\psi^{(2)}(\tfrac14+it)\bigr]
&=8\pi^3\sech^3(2\pi t)-4\pi^3\sech(2\pi t).
\end{align*}
More generally, for $0<c<1$, let $z=(1-c)/2+\iu t$ and
$R_c(t)=\pi\sin(\pi c)/(\cosh(2\pi t)+\cos(\pi c))$.
Reflection and conjugation give $\Real(\psi(z+c)-\psi(z))=R_c(t)$.
Differentiating $m$ times yields
\[
 \Real\!\left(\iu^m[\psi^{(m)}(z+c)-\psi^{(m)}(z)]\right)=R_c^{(m)}(t).
\]
Thus the real part is determined for even $m$ and the imaginary part for odd
$m$. With weight $w=m+1$, these are the real component at odd weight and
imaginary component at even weight. Their complementary components are not
determined by this reflection identity; that leaves open reductions by other
methods. This shifted-difference statement is distinct from individual
polygamma values on $\Real z=1/2$.

\section{Maximal symmetry: vanishing identities and \texorpdfstring{$\Gamma$}{Gamma}-periods}
\label{cm:sec:maximal}

At the two CM points with extra symmetry --- $\tau=i$ (square lattice, $j=1728$) and
$\tau=\rho=e^{i\pi/3}$ (hexagonal lattice, $j=0$) --- classical theory pins every $G_{2k}$
exactly. The square lattice admits no weight-$6$ form ($i\cdot(\mathbb{Z}+\mathbb{Z}i)
=\mathbb{Z}+\mathbb{Z}i$ forces $G_6(i)=i^{-6}G_6(i)=-G_6(i)$), and the hexagonal lattice no
weight-$4$ form; $G_4(i)$ and $G_6(\rho)$ are the lemniscatic and equianharmonic period
powers --- monomials in $\Gamma(1/4)$, resp.\ $\Gamma(1/3)$, by Chowla--Selberg \cite{cm:ChowlaSelberg}.

\subsection{Vanishing identities}

\begin{result}[the CM gems]\label{cm:res:vanishing}
\[
\sum_{n\ge 1}\Bigl[\psi^{(5)}(ni)+\psi^{(5)}(1-ni)\Bigr]=-120\,\zeta(6)
\qquad\text{\dps{residual $1.7\times10^{-49}$}},
\]
\[
\sum_{n\ge 1}\Bigl[\psi^{(3)}(n\rho)+\psi^{(3)}(1-n\rho)\Bigr]=-6\,\zeta(4),
\qquad \rho=e^{i\pi/3}
\qquad\text{\dps{residual $6.8\times10^{-50}$}}.
\]
\end{result}

Each individual bracket is a transcendental-looking complex number (elementary in
$\coth/\operatorname{csch}$ of $\pi n\sqrt3/2$ and the like after the bilateral cotangent
expansion, but nobody would guess the totals); the sums collapse to rational multiples of
$\zeta(6)$, $\zeta(4)$ \emph{because lattice rotation multiplies the corresponding value by a nontrivial root of unity}. Via
\eqref{cm:eq:rowsum} these are exactly the statements $G_6(i)=0$ and $G_4(\rho)=0$.

\subsection{Chowla--Selberg made polygamma-visible}

\begin{result}[period identities]\label{cm:res:periods}
\[
2\zeta(4)+\frac13\sum_{n\ge 1}\Bigl[\psi^{(3)}(ni)+\psi^{(3)}(1-ni)\Bigr]
=\frac{\Gamma(1/4)^8}{960\,\pi^2}
\qquad\text{\dps{residual $2.1\times10^{-50}$}},
\]
\[
2\zeta(6)+\frac1{60}\sum_{n\ge 1}\Bigl[\psi^{(5)}(n\rho)+\psi^{(5)}(1-n\rho)\Bigr]
=\frac{\Gamma(1/3)^{18}}{8960\,\pi^6}
\qquad\text{\dps{residual $2.0\times10^{-60}$}}.
\]
\end{result}

Polygamma values at purely imaginary (resp.\ hexagonal) points, summed over the integer
multiples of the point, know the lemniscate constant and its equianharmonic sibling. The
atoms $\Gamma(1/4)^2/\pi$ and $\Gamma(1/3)$ are precisely the CM-period currency of the
repository's hypergeometric reducer (\texttt{cmPeriodReduce} in \texttt{src/Hypergeometric}),
so these identities tie the polygamma world directly to that corpus.

\subsection{The complete low-weight table}

The symmetry argument and the classical Hurwitz recursion \cite{cm:Hurwitz}
\[
(2n+1)(2n-1)(n-3)\,G_{2n}=3\sum_{m=2}^{n-2}(2m-1)(2n-2m-1)\,G_{2m}G_{2n-2m}
\]
extend both subsections through weight~$12$ with no new constants
\dps{all entries verified, residuals $10^{-60}$--$10^{-263}$}:

\begin{center}
\begin{tabular}{c l l}
\toprule
weight $2k$ & square lattice $\tau=i$ & hexagonal lattice $\tau=\rho$\\
\midrule
$4$ & $\Gamma(1/4)^8/(960\pi^2)$ & $\mathbf{0}$\\
$6$ & $\mathbf{0}$ & $\Gamma(1/3)^{18}/(8960\pi^6)$\\
$8$ & $\tfrac37 G_4(i)^2=3\,\Gamma(1/4)^{16}/(7\cdot 960^2\pi^4)$ & $\mathbf{0}$\\
$10$ & $\mathbf{0}$ & $\mathbf{0}$\\
$12$ & $\tfrac{1}{143}\bigl(18\,G_4(i)^3+25\,G_6(i)^2\bigr)$ type
     & $\tfrac{25}{143}G_6(\rho)^2\propto\Gamma(1/3)^{36}$\\
\bottomrule
\end{tabular}
\end{center}

Every bold zero is a vanishing identity
$\sum_{n\ge1}[\psi^{(w-1)}(n\tau)+\psi^{(w-1)}(1-n\tau)]=-(w-1)!\,\zeta(w)$, and every
non-zero entry is a $\psi^{(w-1)}$-sum evaluation. The pattern is exhaustive: $G_{2k}(i)=0$
iff $2k\not\equiv 0\pmod 4$, $G_{2k}(\rho)=0$ iff $2k\not\equiv 0\pmod 6$, and all non-zero
values are $\mathbb{Q}$-monomials in the two period powers --- so the \emph{entire} Eisenstein
tower over both lattices is polygamma-explicit.

\section{Beyond maximal symmetry: \texorpdfstring{$2i$ and $i\sqrt2$}{2i and i sqrt(2)}}
\label{cm:sec:nextpoints}

The same row-sum at the next two CM points, with closed forms found by PSLQ
\dps{$120$-dps canaries}:

\begin{result}\label{cm:res:2i}
\begin{gather*}
G_4(2i)=\frac{11}{15360}\,\frac{\Gamma(1/4)^8}{\pi^2}
\quad\text{\dps{canary $0$ at $120$ dps}},\\
G_4(i\sqrt2)=\frac{\Gamma(1/8)^4\,\Gamma(3/8)^4}{4608\,\pi^2}
\quad\text{\dps{canary $3.9\times10^{-121}$}}.
\end{gather*}
The weight-6 row at the same points \dps{$160$-dps canaries}:
\[
G_6(2i)=\frac{\Gamma(1/4)^{12}}{2^{14}\cdot 5\cdot\pi^3},\qquad
G_6(i\sqrt2)=\frac{\Gamma(1/8)^6\,\Gamma(3/8)^6}{2^{12}\cdot 3^3\cdot 5\cdot\pi^3}.
\]
\end{result}

The first value confirms the isogeny picture: $2i$ sits over the same field $\mathbb{Q}(i)$,
so only the rational factor changes ($G_4(i)/G_4(2i)=16/11$). The second imports a genuinely
\emph{new period}: discriminant $-8$, field $\mathbb{Q}(\sqrt{-2})$ --- the polygamma sum
$\sum_n[\psi^{(3)}(ni\sqrt2)+\psi^{(3)}(1-ni\sqrt2)]$ evaluates in $\Gamma(1/8)\Gamma(3/8)$.
The fundamental unit $1+\sqrt2$ was offered in the PSLQ basis and declined (exponent $0$).
Together with Section~\ref{cm:sec:maximal} this makes three distinct Chowla--Selberg periods
polygamma-visible.

\section{Class number 2: discriminants \texorpdfstring{$-20$ and $-15$}{-20 and -15}}
\label{cm:sec:h2}

At $\mathbb{Q}(\sqrt{-5})$ (class number $2$) a monomial-only logarithmic basket can miss the value: the algebraic factor in $G_4/\Omega^4$ lives in the Hilbert class field and is not a
monomial. The structure that works is the Hecke/Damerell pair \cite{cm:Damerell} over the two reduced forms
$x^2+5y^2$ ($\tau_1=i\sqrt5$) and $2x^2+2xy+3y^2$ ($\tau_2=(-1+i\sqrt5)/2$).

\begin{result}[disc $-20$]\label{cm:res:d20}
\[
\frac{G_4(\tau_1)}{G_4(\tau_2)}=\frac{29+12\sqrt5}{44}=\frac{(3+2\sqrt5)^2}{44}
\quad\text{\dps{canary $0$ at $160$ dps}},
\]
\[
G_4(\tau_1)\,G_4(\tau_2)
=\frac{11\,\varphi^4\,\Gamma(3/20)^8\,\Gamma(7/20)^8}{2^{14}\cdot3^4\cdot5^4\cdot\pi^4}
\quad\text{\dps{canary $5.3\times10^{-160}$}},
\]
where $\varphi$ is the golden ratio. The ratio is algebraic of degree $2$ with norm
$N(29+12\sqrt5)=121=11^2$; the class product is a period monomial carrying the prime $11$
and a golden unit.
\end{result}

\begin{result}[disc $-15$]\label{cm:res:d15}
With $\tau_1=(1+i\sqrt{15})/2$ for $x^2+xy+4y^2$ and the conjugate pair
$\tau_{2,3}=(\pm1+i\sqrt{15})/4$ representing $2x^2+xy+2y^2$, write
$\widehat G_4=\tfrac12\bigl[G_4(\tau_2)+G_4(\tau_3)\bigr]
=\operatorname{Re}G_4(\tau_2)$ for the class-pair average. Then
\dps{$160$-dps canaries}:
\[
\frac{G_4(\tau_1)}{\widehat G_4}=\frac{2\,(4+\sqrt5)^2}{77},\qquad
G_4(\tau_1)\,\widehat G_4
=\frac{77\,\bigl(\Gamma(\tfrac1{15})\Gamma(\tfrac2{15})\Gamma(\tfrac4{15})
\Gamma(\tfrac8{15})\bigr)^4}{2^{13}\cdot3^6\cdot5^5\cdot\pi^4}.
\]
Unlike disc $-20$, the non-principal value here is genuinely complex (the form
$(2,1,2)$ has $|\tau_2|=1$, where conjugation acts nontrivially on $G_4$); the
class-pair average is the Galois-natural object.
\end{result}

The $\Gamma$-block in Result~\ref{cm:res:d15} is exactly the set of residues
$\{1,2,4,8\}$ with $\chi_{-15}=+1$; the golden unit is declined by PSLQ (exponent $0$); and
the prime content $77=7\cdot11$ of the product matches the ratio's denominator
($N(4+\sqrt5)=11$). The heuristic that emerges, used at every later discriminant:
\emph{run the ratio first; its norm names the primes the class-product basis needs.}

\section{Class number 3: discriminant \texorpdfstring{$-23$}{-23}}\label{cm:sec:h3}

$\mathbb{Q}(\sqrt{-23})$ has $h=3$, with Hilbert class field generated over $\mathbb Q(\sqrt{-23})$ by a root of $x^3-x-1$. The three reduced forms give $\tau_1=(1+i\sqrt{23})/2$ and
$\tau_{2,3}=(\pm1+i\sqrt{23})/4$; $G_4(\tau_1)$ is real and $G_4(\tau_{2,3})$ form a
conjugate pair.

\begin{result}[disc $-23$ class product]\label{cm:res:d23}
\dps{canary $1.4\times10^{-160}$}
\[
G_4(\tau_1)\,G_4(\tau_2)\,G_4(\tau_3)
=\frac{11\cdot17\;P^4}{2^{24}\cdot3^6\cdot23^5\cdot\pi^{16}},\qquad
P=\!\!\prod_{\substack{k\bmod 23\\ (\frac{k}{23})=+1}}\!\!\Gamma\!\Bigl(\frac{k}{23}\Bigr)
\quad\text{(11 factors)}.
\]
\end{result}

\begin{result}[the $\mathfrak{g}_2$ class polynomial from row-sums]\label{cm:res:d23poly}
With $\mathfrak{g}_2=E_4/\eta^8$ phase-corrected by the canonical cube root (the $q^{1/3}$
ambiguity of $\eta^8$ at these points), the three values $\mathfrak{g}_2(\tau_i)$ computed from the
polygamma row-sums are exactly the roots of the classical class polynomial
\[
x^3+155x^2+650x+23375=0
\qquad\text{\dps{coefficients integer to $10^{-80}$}},
\]
whose constant term $23375=5^3\cdot11\cdot17$ forecast the primes $11$, $17$ in the class
product --- the norm heuristic of Section~\ref{cm:sec:h2}, one degree up.
\end{result}

\begin{negative}[weight-4 sums are not Galois-clean]\label{cm:neg:sums}
The dimensionless symmetric functions $e_1^3/e_3$ and $e_2^3/e_3^2$ of the raw $G_4$ values
were not recognized as rational in the source search (height $10^{30}$, $300$ digits). Unlike the value
\emph{product} (a norm) and the weight-$0$ $\mathfrak{g}_2$ invariants, sums of weight-$4$ CM
values across ideal classes are not Galois-stable --- the period correction under
$\operatorname{Gal}(H/K)$ is nontrivial at weight $4$. Class polynomials must be formed from
$\mathfrak{g}_2=E_4/\eta^8$ (or $j$), not from the $G_4$ values themselves.
\end{negative}

\section{Class number 4: discriminant \texorpdfstring{$-39$}{-39} and genus structure}
\label{cm:sec:h4}

The four reduced forms of discriminant $-39$, namely $(1,1,10)$, $(2,\pm1,5)$, $(3,3,4)$,
give $\tau\in\{(1+i\sqrt{39})/2,\ (\pm1+i\sqrt{39})/4,\ (3+i\sqrt{39})/6\}$.

\begin{result}[disc $-39$ class product]\label{cm:res:d39}
\dps{$160$-dps canary}
\[
\prod_{i=1}^{4}G_4(\tau_i)
=\frac{17\cdot23\cdot29\;P^4}{2^{28}\cdot3^{9}\cdot5^{4}\cdot13^{8}\cdot\pi^{16}},\qquad
P=\!\!\prod_{\substack{k\bmod 39\\ \chi_{-39}(k)=+1}}\!\!\Gamma\!\Bigl(\frac{k}{39}\Bigr)
\quad\text{(12 factors)}.
\]
\end{result}

A wrinkle relative to Section~\ref{cm:sec:h3}: $\mathfrak{g}_2$ is \emph{not} a class invariant here
($3\mid 39$, the classical obstruction \cite{cm:Cox}; observed directly --- no cube-root phase choice makes
the $\mathfrak{g}_2$-sums rational). But $j$ always is, and the row-sum values deliver the Hilbert
class polynomial outright.

\begin{result}[Hilbert class polynomial of disc $-39$]\label{cm:res:H39}
\[
H_{-39}(x)=x^4+331531596\,x^3-429878960946\,x^2+109873509788637459\,x
+20919104368024767633
\]
\dps{$e_1,\dots,e_3$ PSLQ-exact integers; $e_4$ integral to $10^{-137}$}.
\end{result}

Worth dwelling on: a quartic with $20$-digit integer coefficients recovered from sums of
polygamma values at quadratic-irrational complex points. The polygamma row-sum is a
perfectly serviceable special-function front door to the classical CM tower.

\begin{result}[genus structure inside the quartet]\label{cm:res:genus}
The class group is $\mathbb{Z}/4$, with the ambiguous form $(3,3,4)$ as the order-$2$ class.
The raw cross-class ratio $v_1/v_4$ is a quartic candidate; no degree-two polynomial was found up to height $10^{12}$, but the genus-wise ratio --- principal genus $\{1,A^2\}$
over $\{A,A^3\}$ --- is exactly quadratic, in the real part of the genus field
$\mathbb{Q}(\sqrt{-3},\sqrt{13})$ as genus theory demands
\dps{$160$-dps canary; $\sqrt3$, $\sqrt{39}$ declined by PSLQ}:
\[
\frac{G_4(\tau_1)\,G_4(\tau_4)}{G_4(\tau_2)\,G_4(\tau_3)}
=\frac{103299+4440\sqrt{13}}{181424},\qquad
N=\Bigl(\frac{9}{16}\Bigr)^{\!2},\qquad 181424=16\cdot17\cdot23\cdot29 .
\]
The denominator reuses the class-product primes $17$, $23$, $29$.
\end{result}

\section{Class number 5: discriminant \texorpdfstring{$-47$}{-47}}\label{cm:sec:h5}

The five reduced forms $(1,1,12)$, $(2,\pm1,6)$, $(3,\pm1,4)$ give the quintet
$\tau\in\{(1+i\sqrt{47})/2,\ (\pm1+i\sqrt{47})/4,\ (\pm1+i\sqrt{47})/6\}$.

\begin{result}[disc $-47$ class product]\label{cm:res:d47}
\dps{$200$-dps canary}
\[
\prod_{i=1}^{5}G_4(\tau_i)
=\frac{11^2\cdot23\cdot29\;P^4}{2^{52}\cdot3^{6}\cdot47^{9}\cdot\pi^{36}},\qquad
P=\!\!\prod_{\substack{k\bmod 47\\ (\frac{k}{47})=+1}}\!\!\Gamma\!\Bigl(\frac{k}{47}\Bigr)
\quad\text{(23 factors)}.
\]
\end{result}

\begin{result}[Hilbert class polynomial of disc $-47$]\label{cm:res:H47}
The row-sum $j$-values deliver the quintic Hilbert class polynomial
\dps{all symmetric functions integral to $10^{-129}$ or better}:
\begin{align*}
H_{-47}(x)=x^5&+2257834125\,x^4-9987963828125\,x^3+5115161850595703125\,x^2\\
&-14982472850828613281250\,x+16042929600623870849609375 .
\end{align*}
\end{result}

The CM tower is now polygamma-explicit through class number $5$ with one uniform recipe:
\emph{row-sum $\to$ class-product log-PSLQ (norm-forecast primes) $\to$
$j$-symmetric-function integrality.}

\section{The opaque components}\label{cm:sec:opaque}

The components \emph{not} fixed by reflection-conjugation or by the even-row mechanism ---
$\operatorname{Im}\psi^{(1)}(i)$, $\operatorname{Re}\psi(i)$ beyond $\gamma$,
$\operatorname{Im}\psi^{(1)}\bigl(\tfrac{1+i}{2}\bigr)$ --- are \emph{odd-row} lattice data:
for example
\[
\operatorname{Im}\psi^{(1)}(i)=-2\sum_{k\ge1}\frac{k}{(k^2+1)^2}
\]
is an odd-power analogue of the $\coth$ family, which the bilateral (even) cotangent
expansion cannot reach. These are Eichler-integral/quasi-modular territory --- the
$E_2$-flavoured layer; compare $\sum_{k\ge1}k/(e^{2\pi k}-1)=\tfrac1{24}-\tfrac1{8\pi}$,
which is $E_2(i)=3/\pi$ in Lambert-series clothing.

\begin{negative}\label{cm:neg:opaque}
A PSLQ basket for these constants against products drawn from
\[
\{1,\ \pi,\ \coth\pi,\ \pi^2/\sinh^2\pi,\ \pi\coth\pi,\ \dots\}
\]
found nothing at small coefficient heights --- consistent with genuinely new constants.
Recorded as an honest negative; the Lambert-series face ($\sum\sigma_5(n)q^n$ at
$q=e^{-2\pi}$ and $q=-e^{-\pi\sqrt3}$) is the natural next probe surface.
\end{negative}

\section{Summation order, differentiation and certified tails}
Absolute convergence made the preceding weights straightforward. Weight
two needs a specified order of summation. Once that boundary is handled,
Ramanujan's differential algebra turns the existing CM period seeds into
all derivative evaluations, with explicit exponential tail estimates.

\section{Regularized modular rows and all derivative evaluations}
\label{modular:sec:modular}

The modular draft obtains even-weight Eisenstein series by pairing
polygamma values along lattice rows. Its weight-two question requires
care: the paired rows converge rapidly, although the ungrouped lattice
sum is not absolutely convergent. We first prove an explicit bound for
every differentiated row. We then determine the change of summation
order directly, before using any modular transformation law. Finally,
Ramanujan's differential system evaluates all the resulting derivative
rows at the square and hexagonal CM points.

Throughout this section, let
\[
\tau\in\mathbb H:=\{z\in\C:\Imag z>0\},\qquad
q=e^{2\pi\iu\tau},\qquad r=|q|=e^{-2\pi\Imag\tau}<1,
\qquad \rho=\tau_{\mathrm{hex}}=\frac{1+\iu\sqrt3}{2}.
\]
For even \(w\ge2\) and an integer \(j\ge0\), put
\begin{equation}
\mathcal T_{w,j}(\tau)
=\sum_{n\ge1}n^j\left[
 \psi^{(w-1+j)}(n\tau)
 +(-1)^j\psi^{(w-1+j)}(1-n\tau)\right].
\label{modular:eq:modular-T-definition}
\end{equation}
Convergence will follow from the next theorem. The sign on the reflected
term is part of the definition; differentiating that term changes its
sign at every step.

\subsection{Rational row formulas and explicit tail bounds}

The negative-index polylogarithms appearing below are elementary
rational functions. For \(d\ge0\), define
\[
R_d(t)=\Li_{-d}(t),\qquad
R_0(t)=\frac{t}{1-t},\qquad
R_{d+1}(t)=t\frac{d}{dt}R_d(t).
\]
These relations are identities of rational functions, since on
\(|t|<1\) they follow by differentiating
\(R_d(t)=\sum_{\ell\ge1}\ell^d t^\ell\). In particular,
\[
R_1(t)=\frac{t}{(1-t)^2},\qquad
R_2(t)=\frac{t(1+t)}{(1-t)^3},\qquad
R_3(t)=\frac{t(1+4t+t^2)}{(1-t)^4}.
\]

\begin{theorem}[Rational rows and an explicit remainder]
\label{modular:thm:modular-rows-tail}
For \(w,j,\tau\) as above,
\begin{align}
&\psi^{(w-1+j)}(n\tau)
 +(-1)^j\psi^{(w-1+j)}(1-n\tau)
 =(2\pi\iu)^{w+j}\Li_{1-w-j}(q^n),
\label{modular:eq:modular-single-row}\\
&\mathcal T_{w,j}(\tau)
 =(2\pi\iu)^{w+j}\sum_{n\ge1}n^j\Li_{1-w-j}(q^n).
\label{modular:eq:modular-rational-rows}
\end{align}
The series converges normally on every compact subset of \(\mathbb H\).
If \(\mathcal T_{w,j}^{[N]}\) is its sum through \(n=N\), where
\(N\ge0\), then
\begin{equation}
\left|\mathcal T_{w,j}(\tau)-\mathcal T_{w,j}^{[N]}(\tau)\right|
\le (2\pi)^{w+j}(N+1)^j
 \frac{\Li_{-j}(r)}{r}\,
 \Li_{1-w-j}\!\left(r^{N+1}\right).
\label{modular:eq:modular-tail-bound}
\end{equation}
Thus the bound contains only rational functions of \(r\) and
\(r^{N+1}\), together with an explicit power of \(\pi\).
\end{theorem}

\begin{proof}
For \(\Imag z>0\), splitting the absolutely convergent bilateral
power sum into its nonnegative and negative indices gives
\begin{equation}
\sum_{m\in\Z}(m+z)^{-w}
=\frac{\psi^{(w-1)}(z)+\psi^{(w-1)}(1-z)}{(w-1)!}.
\label{modular:eq:modular-bilateral-polygamma}
\end{equation}
Here evenness of \(w\) is essential. The classical cotangent expansion
and its partial-fraction interpretation~\cite{hstruct:DLMF} read
\[
\pi\cot(\pi z)=-\pi\iu
 -2\pi\iu\sum_{\ell\ge1}e^{2\pi\iu\ell z},\qquad
\pi\cot(\pi z)=\lim_{M\to\infty}\sum_{m=-M}^{M}\frac1{m+z}.
\]
Differentiate \(w-1\) times. On compact subsets of the upper
half-plane the exponential series and its derivatives converge
uniformly. After the first derivative, the bilateral series and all
its subsequent derivatives also converge normally. Since
\((-1)^{w-1}=-1\), the two minus signs cancel and yield
\[
\psi^{(w-1)}(z)+\psi^{(w-1)}(1-z)
=(2\pi\iu)^w\sum_{\ell\ge1}\ell^{w-1}e^{2\pi\iu\ell z}.
\]
Differentiating another \(j\) times with respect to \(z\) proves
\eqref{modular:eq:modular-single-row}. Multiplication by \(n^j\) gives the
summand in \eqref{modular:eq:modular-rational-rows}.

To prove the remainder bound, expand the rational polylogarithm in
its nonnegative majorizing series. The absolute value of the omitted
rows is at most
\begin{equation}
(2\pi)^{w+j}\sum_{\ell\ge1}\ell^{w+j-1}
                  \sum_{n>N}n^j r^{n\ell}.
\label{modular:eq:modular-positive-majorant}
\end{equation}
Write \(n=N+1+a\). For \(a\ge0\) and \(\ell\ge1\),
\[
n^j\le(N+1)^j(a+1)^j,\qquad r^{a\ell}\le r^a.
\]
Consequently
\[
\sum_{n>N}n^j r^{n\ell}
\le (N+1)^j r^{(N+1)\ell}
       \sum_{a\ge0}(a+1)^j r^a
=(N+1)^j r^{(N+1)\ell}\frac{\Li_{-j}(r)}{r}.
\]
Substitution in \eqref{modular:eq:modular-positive-majorant} proves
\eqref{modular:eq:modular-tail-bound}. In particular, the majorant is finite.
On a compact subset of \(\mathbb H\), one has \(r\le r_0<1\), and
each of its power series with nonnegative coefficients is bounded by
its value at \(r_0\). This proves normal convergence, for every fixed
\(w,j\), and justifies all subsequent termwise differentiations.
\end{proof}

For \(j=0\), the tail formula simplifies to
\[
\left|\mathcal T_{w,0}-\mathcal T_{w,0}^{[N]}\right|
\le\frac{(2\pi)^w}{1-r}\Li_{1-w}\!\left(r^{N+1}\right).
\]
The case \(w=2\), which will be particularly useful, is
\begin{equation}
\left|\mathcal T_{2,0}-\mathcal T_{2,0}^{[N]}\right|
\le\frac{4\pi^2r^{N+1}}
 {(1-r)(1-r^{N+1})^2}.
\label{modular:eq:modular-weight-two-tail}
\end{equation}
Theorem~\ref{modular:thm:modular-rows-tail} also proves
\begin{equation}
\frac{d^j}{d\tau^j}\mathcal T_{w,0}(\tau)
=\mathcal T_{w,j}(\tau).
\label{modular:eq:modular-termwise-derivatives}
\end{equation}
Computationally, \eqref{modular:eq:modular-single-row} is preferable to
evaluating its two polygamma terms separately. Their algebraic
asymptotic contributions cancel, leaving the exponentially small
paired row. The rational formula performs that cancellation exactly
before numerical evaluation.

\subsection{The weight-two ordering correction from telescoping}

For \(w\ge4\), the ordinary absolutely convergent lattice series is
\[
G_w(\tau)=\sum_{(m,n)\in\Z^2\setminus\{(0,0)\}}
                  (m+n\tau)^{-w}
=2\zeta(w)+\frac{2}{(w-1)!}\mathcal T_{w,0}(\tau).
\]
For \(w=2\), specify the order explicitly:
\begin{align}
G_2(\tau)&=\sum_{n\in\Z}
        \left(\sum_{m\in\Z}'(m+n\tau)^{-2}\right)
 =\frac{\pi^2}{3}+2\mathcal T_{2,0}(\tau),
\label{modular:eq:modular-G2-definition}\\
C_2(\tau)&=\sum_{m\in\Z}
        \left(\sum_{n\in\Z}'(m+n\tau)^{-2}\right).
\label{modular:eq:modular-C2-definition}
\end{align}
The prime omits only \((m,n)=(0,0)\). In the first expression, each
inner row and the outer series of completed rows converge absolutely.
The reverse expression exists as well: changing the names and signs
of the one-dimensional indices gives
\begin{equation}
C_2(\tau)=\tau^{-2}G_2(-1/\tau).
\label{modular:eq:modular-C2-coordinate}
\end{equation}
This is a coordinate identity for the prescribed iterated sums, not
an assumption of modularity. The series of absolute values over all
lattice points diverges, so an interchange of the two orders still
requires a calculation.

\begin{theorem}[Summation-order anomaly]
\label{modular:thm:modular-ordering}
For every \(\tau\in\mathbb H\),
\begin{equation}
C_2(\tau)=G_2(\tau)-\frac{2\pi\iu}{\tau}.
\label{modular:eq:modular-anomaly}
\end{equation}
Consequently
\begin{equation}
G_2(-1/\tau)=\tau^2G_2(\tau)-2\pi\iu\tau,
\qquad G_2(\tau+1)=G_2(\tau).
\label{modular:eq:modular-G2-transform}
\end{equation}
\end{theorem}

\begin{proof}
We prove the difference of the two orders by replacing their common
part with an absolutely convergent series. Put
\[
h(z)=\frac1z-\frac1{z+1},\qquad
f(z)=\frac1{z^2}-h(z)=\frac1{z^2(z+1)},
\]
and define
\[
F(\tau)=\sum_{\substack{n\in\Z\setminus\{0\}\\m\in\Z}}
                    f(m+n\tau).
\]
No pole occurs in this sum. The real-linear map
\((m,n)\mapsto m+n\tau\) is invertible, so for some \(c_\tau>0\),
\(|m+n\tau|\ge c_\tau\sqrt{m^2+n^2}\).
Outside a finite set, \(|f(z)|\le2|z|^{-3}\); hence \(F\) converges
absolutely. It can therefore be summed in either order.

For each fixed \(n\ne0\), the series \(\sum_m h(m+n\tau)\) is
absolutely convergent and its symmetric partial sum telescopes:
\[
\sum_{m=-M}^{M}h(m+n\tau)
=\frac1{-M+n\tau}-\frac1{M+1+n\tau}\longrightarrow0.
\]
Separating the row \(n=0\) proves
\begin{equation}
G_2(\tau)=\frac{\pi^2}{3}+F(\tau).
\label{modular:eq:modular-G2-absolute-part}
\end{equation}

For the reverse order, define for integer \(m\)
\[
A(m)=\lim_{K\to\infty}
              \sum_{0<|n|\le K}\frac1{m+n\tau}.
\]
This symmetric limit exists: pairing \(n\) and \(-n\) gives an
\(O(n^{-2})\) summand. The cotangent partial-fraction formula gives
\begin{equation}
A(m)=\frac{\pi}{\tau}\cot\!\left(\frac{\pi m}{\tau}\right)-\frac1m
\quad(m\ne0),\qquad A(0)=0.
\label{modular:eq:modular-A-cotangent}
\end{equation}
The value at zero is also the removable limit of the expression on
the right. For each fixed \(m\), the \(n\)-series of \(h(m+n\tau)\)
is absolutely convergent. Taking its symmetric partial sums shows
\[
\sum_{n\ne0}h(m+n\tau)=A(m)-A(m+1).
\]
It follows that
\begin{align*}
C_2(\tau)-\frac{\pi^2}{3}-F(\tau)
&=\lim_{M\to\infty}\sum_{m=-M}^{M}\sum_{n\ne0}h(m+n\tau)\\
&=\lim_{M\to\infty}\bigl(A(-M)-A(M+1)\bigr).
\end{align*}
The extraction of \(F\) here is legitimate by its absolute
convergence; no rearrangement of the original double power sum has
been made.

Since \(\Imag(1/\tau)<0\),
\[
\cot(\pi m/\tau)\longrightarrow\iu\quad(m\to+\infty),
\qquad
\cot(\pi m/\tau)\longrightarrow-\iu\quad(m\to-\infty).
\]
Equation~\eqref{modular:eq:modular-A-cotangent} therefore gives
\(A(-M)-A(M+1)\to-2\pi\iu/\tau\). Combining this with
\eqref{modular:eq:modular-G2-absolute-part} proves
\eqref{modular:eq:modular-anomaly}. Equation~\eqref{modular:eq:modular-C2-coordinate}
now gives the first transformation in
\eqref{modular:eq:modular-G2-transform}. Periodicity follows directly from
the \(q\)-series in \eqref{modular:eq:modular-rational-rows}.
\end{proof}

The anomaly is classical; Romik and Scherer discuss it together with
more general shape-dependent summation procedures~\cite{modular:RomikScherer}.
The proof above fixes both its sign and the normalization needed for
the draft's polygamma sums.

\begin{corollary}[The missing weight-two CM rows]
\label{modular:cor:modular-weight-two-CM}
One has
\begin{align}
G_2(\iu)&=\pi,&G_2(\rho)&=\frac{2\pi}{\sqrt3},
\label{modular:eq:modular-G2-CM}\\
\mathcal T_{2,0}(\iu)&=\frac\pi2-\frac{\pi^2}{6},&
\mathcal T_{2,0}(\rho)&=\frac\pi{\sqrt3}-\frac{\pi^2}{6}.
\label{modular:eq:modular-trigamma-CM}
\end{align}
\end{corollary}
\begin{proof}
At \(\tau=\iu\), the first transformation in
\eqref{modular:eq:modular-G2-transform} gives
\(G_2(\iu)=-G_2(\iu)+2\pi\).
For \(\rho\), use \(-1/\rho=\rho-1\) and periodicity. The same
transformation becomes
\[
(\rho^2-1)G_2(\rho)=2\pi\iu\rho.
\]
Since \(\rho^2-1=\iu\sqrt3\rho\), this proves the second value in
\eqref{modular:eq:modular-G2-CM}. Substitution into
\eqref{modular:eq:modular-G2-definition} proves
\eqref{modular:eq:modular-trigamma-CM}.
\end{proof}

\subsection{CM seeds and the complete differential tower}

Normalize the Eisenstein series by
\[
E_2=\frac{3G_2}{\pi^2},\qquad
E_4=\frac{G_4}{2\zeta(4)},\qquad
E_6=\frac{G_6}{2\zeta(6)}.
\]
The seed values, in precisely this normalization, are
\begin{align}
(E_2(\iu),E_4(\iu),E_6(\iu))
&=\left(\frac3\pi,
 \frac{3\Gamma(1/4)^8}{64\pi^6},0\right),
\label{modular:eq:modular-square-seeds}\\
(E_2(\rho),E_4(\rho),E_6(\rho))
&=\left(\frac{2\sqrt3}\pi,0,
 \frac{27\Gamma(1/3)^{18}}{512\pi^{12}}\right).
\label{modular:eq:modular-hex-seeds}
\end{align}
We have already proved the weight-two entries. The zero entries follow
by rotating the absolutely convergent lattices. Multiplication by
\(\iu\) fixes the square lattice and sends its weight-six sum to its
negative; multiplication by \(\rho\) fixes the hexagonal lattice and
multiplies its weight-four sum by \(\rho^{-4}\ne1\).

For completeness, the classical period integrals fix the nonzero
entries without numerical identification. The standard inverse
Weierstrass integral~\cite{hstruct:DLMF}, with
\(g_2=60G_4\) and \(g_3=140G_6\), gives positive real periods
\begin{align*}
p_{\mathrm{sq}}
&=2\int_1^\infty\frac{dx}{\sqrt{4x^3-4x}}
 =2\int_0^1\frac{du}{\sqrt{1-u^4}}
 =\frac12 B\!\left(\frac14,\frac12\right)
 =\frac{\Gamma(1/4)^2}{2\sqrt{2\pi}},\\
p_{\mathrm{hex}}
&=2\int_1^\infty\frac{dx}{\sqrt{4x^3-4}}
 =\frac13 B\!\left(\frac16,\frac12\right)
 =\frac{\Gamma(1/3)^3}{2^{4/3}\pi}.
\end{align*}
The corresponding lattices are
\(p_{\mathrm{sq}}(\Z+\iu\Z)\) for \((g_2,g_3)=(4,0)\), and
\(p_{\mathrm{hex}}(\Z+\rho\Z)\) for \((g_2,g_3)=(0,4)\).
Their square and equianharmonic shapes follow from the rotations of
the respective cubics; the real period integral fixes the scale and
orientation. The first substitution above is \(x=u^{-2}\); in the
second beta integral use \(t=x^{-3}\). Euler's reflection and
duplication formulas give the displayed gamma evaluations.
Scaling the lattice invariants consequently yields
\[
G_4(\iu)=\frac{p_{\mathrm{sq}}^4}{15}
 =\frac{\Gamma(1/4)^8}{960\pi^2},\qquad
G_6(\rho)=\frac{p_{\mathrm{hex}}^6}{35}
 =\frac{\Gamma(1/3)^{18}}{8960\pi^6}.
\]
Using \(2\zeta(4)=\pi^4/45\) and
\(2\zeta(6)=2\pi^6/945\) proves the remaining entries of
\eqref{modular:eq:modular-square-seeds}--\eqref{modular:eq:modular-hex-seeds}.

Ramanujan's classical differential identities~\cite{modular:Ramanujan} are
\begin{equation}
\theta E_2=\frac{E_2^2-E_4}{12},\qquad
\theta E_4=\frac{E_2E_4-E_6}{3},\qquad
\theta E_6=\frac{E_2E_6-E_4^2}{2},\qquad
\theta=\frac1{2\pi\iu}\frac{d}{d\tau}.
\label{modular:eq:modular-Ramanujan-system}
\end{equation}
Encode these identities by the rational derivation
\begin{equation}
\mathcal V
=\frac{X^2-Y}{12}\frac{\partial}{\partial X}
 +\frac{XY-Z}{3}\frac{\partial}{\partial Y}
 +\frac{XZ-Y^2}{2}\frac{\partial}{\partial Z},
\qquad P_0=X,\quad P_{j+1}=\mathcal V P_j.
\label{modular:eq:modular-P-recurrence}
\end{equation}
For example,
\begin{align*}
P_1&=\frac{X^2-Y}{12},\\
P_2&=\frac{X^3-3XY+2Z}{72},\\
P_3&=\frac{X^4-6X^2Y+8XZ-3Y^2}{288}.
\end{align*}

\begin{theorem}[All differentiated weight-two rows]
\label{modular:thm:modular-all-jets}
Set \(\mathcal S_j(\tau)=\mathcal T_{2,j}(\tau)\). For every
\(\tau\in\mathbb H\),
\begin{align}
\mathcal S_0(\tau)&=\frac{\pi^2}{6}\bigl(E_2(\tau)-1\bigr),
\label{modular:eq:modular-S0}\\
\mathcal S_j(\tau)&=\frac{\pi^2}{6}(2\pi\iu)^j
 P_j(E_2(\tau),E_4(\tau),E_6(\tau)),\qquad j\ge1.
\label{modular:eq:modular-all-jets}
\end{align}
In particular, substitution of
\eqref{modular:eq:modular-square-seeds}--\eqref{modular:eq:modular-hex-seeds}
evaluates every \(\mathcal S_j(\iu)\) and
\(\mathcal S_j(\rho)\) by a finite polynomial calculation over
\(\Q\) and the indicated gamma periods.
\end{theorem}

\begin{proof}
Equation~\eqref{modular:eq:modular-G2-definition} and the definition of
\(E_2\) give \eqref{modular:eq:modular-S0}. By the chain rule and
\eqref{modular:eq:modular-Ramanujan-system}, every polynomial
\(P\in\Q[X,Y,Z]\) satisfies
\[
\theta\bigl(P(E_2,E_4,E_6)\bigr)
=(\mathcal V P)(E_2,E_4,E_6).
\]
Induction from \(P_0=X\) therefore gives
\(\theta^jE_2=P_j(E_2,E_4,E_6)\).
Differentiate \eqref{modular:eq:modular-S0} \(j\ge1\) times and use
\eqref{modular:eq:modular-termwise-derivatives}. The constant term disappears,
and \(d/d\tau=2\pi\iu\theta\) contributes the factor
\((2\pi\iu)^j\). This proves \eqref{modular:eq:modular-all-jets}; the CM
specializations use the established seed values.
\end{proof}

For clarity, the first two differentiated evaluations are
\begin{align}
\sum_{n\ge1}n\left[\psi''(n\iu)-\psi''(1-n\iu)\right]
 &=\frac{\iu\pi}{4}
   -\frac{\iu\Gamma(1/4)^8}{768\pi^3},
\label{modular:eq:modular-S1-square}\\
\sum_{n\ge1}n\left[\psi''(n\rho)-\psi''(1-n\rho)\right]
 &=\frac{\iu\pi}{3},
\label{modular:eq:modular-S1-hex}\\
\sum_{n\ge1}n^2\left[\psi'''(n\iu)+\psi'''(1-n\iu)\right]
 &=-\frac\pi4+\frac{\Gamma(1/4)^8}{256\pi^3},
\label{modular:eq:modular-S2-square}\\
\sum_{n\ge1}n^2\left[\psi'''(n\rho)+\psi'''(1-n\rho)\right]
 &=-\frac{2\sqrt3\pi}{9}
   -\frac{\Gamma(1/3)^{18}}{1024\pi^8}.
\label{modular:eq:modular-S2-hex}
\end{align}
These follow by substituting the seed values into \(P_1,P_2\), so
their derivation contains no integer-relation search.

\begin{corollary}[Polynomial degree bounds for the CM row formulas]
\label{modular:cor:modular-CM-degree}
Put
\(A=\Gamma(1/4)^8/\pi^4\) and
\(B=\Gamma(1/3)^{18}/\pi^9\). For each \(j\ge1\), there are
polynomials \(U_j\in\Q[t]\) and
\(V_j\in\Q(\sqrt3)[t]\) such that
\[
\frac{\mathcal S_j(\iu)}{\iu^j\pi}=U_j(A),\qquad
\frac{\mathcal S_j(\rho)}{\iu^j\pi}=V_j(B),
\]
with
\[
\deg U_j\le\left\lfloor\frac{j+1}{2}\right\rfloor,
\qquad
\deg V_j\le\left\lfloor\frac{j+1}{3}\right\rfloor.
\]
\end{corollary}
\begin{proof}
Give \(X,Y,Z\) weights \(2,4,6\). Each term of
\(\mathcal V\) raises weight by two, so \(P_j\) is homogeneous
of weight \(2j+2\). At the square point, the surviving monomials
\(X^aY^b\) have \(a+2b=j+1\). On inserting the seeds into
\eqref{modular:eq:modular-all-jets} and dividing by \(\iu^j\pi\), each
becomes a rational multiple of \(A^b\). At the hexagonal point,
the surviving monomials \(X^aZ^c\) satisfy \(a+3c=j+1\), and
the same substitution produces an element of
\(\Q(\sqrt3)B^c\). The stated degree bounds follow.
\end{proof}

\subsection{What this resolves and how it is verified}

Corollary~\ref{modular:cor:modular-weight-two-CM} resolves the weight-two
row-sum question posed in the modular draft. Theorem~\ref{modular:thm:modular-all-jets}
extends it to all polynomially weighted derivative rows, with the
explicit error estimate of Theorem~\ref{modular:thm:modular-rows-tail}.
These are evaluations of specified infinite aggregates. They do not
evaluate an isolated component such as \(\Imag\psi'(\iu)\), nor
do they prove that such a component is an Eichler integral. Any proposed
identification of that isolated constant needs its own exact transform
or functional equation.

The accompanying \path{code/modular_rows.py} constructs the rational
functions \(R_d\) and the polynomials \(P_j\) by exact arithmetic.
It checks individual rows against a separate complex-polygamma
implementation, compares the CM formulas for \(j=0,\ldots,4\) with
their rapidly convergent row sums, and checks representative tails.
The output \path{data/modular_rows.json} records the exact coefficients,
analytic truncation bounds, and numerical residuals. Its decimal
calculations use ordinary high-precision arithmetic; they are numerical
checks, not outward-rounded interval certificates. A machine enclosure
can be obtained by evaluating the rational formulas and
\eqref{modular:eq:modular-tail-bound} with directed interval arithmetic. The
identities themselves are proved above and do not depend on the
numerical checks.
